\clearpage\corrigesChapitre{19}{Circuits et lois électriques}
\corrigesSouples
\begin{blocLivre}{cor-19-1}\begin{corrige}{1}{Choisir le branchement}
\schema[corrige]{mesures-circuit}
\begin{questions}
\item Voltmètre en dérivation sur la résistance ; ampèremètre en série avec elle.
\item Le brancher directement entre les bornes du générateur : il pourrait créer un court-circuit.
\end{questions}
\end{corrige}\end{blocLivre}
\begin{blocLivre}{cor-19-2}\begin{corrige}{2}{Convertir une intensité}
\begin{questions}
\item $I=\qty{0,0250}{A}$.
\item $U=RI=200\times0,0250=\qty{5,00}{V}$.
\end{questions}
\end{corrige}\end{blocLivre}
\begin{blocLivre}{cor-19-3}\begin{corrige}{3}{Une dérivation}
\schema[corrige]{noeud}
\begin{questions}
\item $I_2=0,30-0,12=\qty{0,18}{A}$.
\item Loi des nœuds : les intensités entrantes et sortantes s’équilibrent.
\end{questions}
\end{corrige}\end{blocLivre}
\begin{blocLivre}{cor-19-4}\begin{corrige}{4}{Une maille simple}
\begin{questions}
\item $U_2=9,0-3,5=\qty{5,5}{V}$.
\item Non. En série, les tensions s’ajoutent ; leurs valeurs dépendent des dipôles.
\end{questions}
\end{corrige}\end{blocLivre}
\begin{blocLivre}{cor-19-5}\begin{corrige}{5}{Une caractéristique ohmique}
\schema[corrige]{ohm}
\begin{questions}
\item Par exemple $R=4,0/0,020=\qty{200}{\ohm}$. Il faut convertir les mA en A.
\item $U=\qty{4,0}{V}$.
\end{questions}
\end{corrige}\end{blocLivre}
\begin{blocLivre}{cor-19-6}\begin{corrige}{6}{Résistance inconnue}
\begin{questions}
\item $R=6,0/0,015=\qty{400}{\ohm}$.
\item $I=3,0/400=\qty{0,0075}{A}=\qty{7,5}{mA}$.
\end{questions}
\end{corrige}\end{blocLivre}
\begin{blocLivre}{cor-19-7}\begin{corrige}{7}{Deux branches parallèles}
\schema[corrige]{parallele}
\begin{questions}
\item $I_1=6,0/100=60$ mA et $I_2=6,0/200=30$ mA.
\item $I=I_1+I_2=\qty{90}{mA}$.
\end{questions}
\end{corrige}\end{blocLivre}
\begin{blocLivre}{cor-19-8}\begin{corrige}{8}{Inverser les bornes}
\begin{questions}
\item $U_{BA}=-U_{AB}=\qty{-4,2}{V}$.
\item Non : il traduit l’inversion du choix des bornes de mesure.
\end{questions}
\end{corrige}\end{blocLivre}
\begin{blocLivre}{cor-19-9}\begin{corrige}{9}{Deux résistances en série}
\schema[corrige]{serie}
\begin{questions}
\item $9,0=R_1I+R_2I=(R_1+R_2)I$.
\item $I=9,0/300=\qty{0,030}{A}$ ; $U_1=3,0$ V et $U_2=6,0$ V.
\end{questions}
\end{corrige}\end{blocLivre}
\begin{blocLivre}{cor-19-10}\begin{corrige}{10}{Un courant non consommé}
\begin{questions}
\item 20 mA : l’intensité est la même dans tous les dipôles de cette branche série.
\item La charge ne disparaît pas ; la résistance est le siège de transferts d’énergie, ce qui est une autre description.
\end{questions}
\end{corrige}\end{blocLivre}
\begin{blocLivre}{cor-19-11}\begin{corrige}{11}{Un dipôle non ohmique}
\begin{questions}
\item 100 ; 140 ; environ 183 ohms : les quotients ne sont pas constants.
\item Non sur cette gamme : la tension n’est pas proportionnelle à l’intensité.
\end{questions}
\end{corrige}\end{blocLivre}
\begin{blocLivre}{cor-19-12}\begin{corrige}{12}{Une modélisation numérique}
\demarche{Lire les unités de la pente puis comparer le décalage aux écarts de mesure.}
\begin{questions}
\item Une résistance d’environ 198 ohms.
\item Non : 0,02 V est de l’ordre des écarts de mesure. Examiner les données et l’étalonnage avant de rejeter le passage par l’origine.
\end{questions}
\retiens{Une faible ordonnée à l’origine n’invalide pas automatiquement le modèle ohmique.}
\end{corrige}\end{blocLivre}

